\chapter{Literature Review}

This chapter reviews the theoretical and technical foundations that underpin
the design decisions made in this project. It covers seven domains:
e-learning in Vietnamese higher education; the H5P interactive content
framework; interactive video and its pedagogical basis; user experience
design and cognitive load theory; artificial intelligence for educational
content generation; cloud-native web architecture; and LMS integration
standards. The chapter concludes by surveying related work and identifying
the research gaps this project addresses.

% ─────────────────────────────────────────────────────────────────────────────
\section{E-Learning in Vietnamese Higher Education}

\subsection{National Context and Policy}

Vietnam's transition to a knowledge economy has made the quality of higher
education a strategic priority. The National Education Development Strategy
2021--2030 mandates that at least 30\% of university-level programmes be
partially or fully delivered through digital platforms by
2025~\cite{vietnam_strategy2021}. VNU-HCM has been the most active
implementer of this policy within the university sector. Its MOOC platform
(\url{https://courses.vnuhcm.edu.vn}) hosts hundreds of university-level
courses, reaching learners not only in Vietnam but across
Southeast Asia~\cite{vnuhcm_mooc}.

However, policy-driven digital transformation often outpaces the readiness
of the human infrastructure it depends on. Bates argues that sustainable
digital education transformation requires investing not only in
infrastructure but in educator capability---the ability of faculty to design
and produce high-quality digital learning materials~\cite{bates2019}.
Beetham and Sharpe extend this argument, proposing that institutions should
provide educators with design frameworks and tools that make digital pedagogy
accessible without requiring deep technical expertise~\cite{beetham2019}.

\subsection{Adoption Barriers Among Vietnamese Faculty}

Research on e-learning quality in developing-world universities, including
Vietnamese institutions, consistently identifies usability as the primary
adoption barrier~\cite{hadullo2017}. Hadullo et al.\ found that in
developing-country university contexts---with structural similarities to
Vietnamese institutions in terms of faculty digital literacy profiles---poor
tool usability accounted for 41\% of the variance in self-reported e-learning
quality. Faculty members reported that they would invest time in content
creation if and only if the authoring process did not itself require
specialist technical knowledge.

The Technology Acceptance Model (TAM) proposed by Davis provides a formal
framework for understanding these dynamics~\cite{davis1989}. TAM posits that
both perceived usefulness (the degree to which a tool is believed to enhance
performance) and perceived ease of use (the degree to which using a tool is
believed to be free of effort) are necessary conditions for adoption. In
educational technology contexts, perceived ease of use often carries more
weight than perceived usefulness, because educators already believe that
interactive content is beneficial---the question is whether they can
practically produce it.

Blended learning research by Garrison and Vaughan reinforces this picture,
showing that faculty adoption of digital tools for blended course delivery
is most strongly predicted by the simplicity of the first authoring
experience~\cite{garrison2008}. Instructors who fail to complete a task
successfully in their first session rarely attempt the task again.

% ─────────────────────────────────────────────────────────────────────────────
\section{The H5P Interactive Content Framework}

\subsection{Architecture and Content Types}

H5P (HTML5 Package) is an open-source framework initially released in 2013
by Joubel AS (later the H5P Group)~\cite{h5p2023}. Its architecture is
built around self-contained content packages: each \texttt{.h5p} file is a
ZIP archive containing all JavaScript and CSS library files required to
render the content, the content parameters as a JSON document, and any
embedded media. This self-contained structure makes H5P content portable
across any web environment that can serve a static file.

The H5P library ecosystem currently encompasses more than 50 content types.
Relevant to this project are the six types listed in Table~\ref{tab:h5p_types}.

\begin{table}[H]
  \centering
  \caption{H5P content types supported by the platform}
  \label{tab:h5p_types}
  \begin{tabularx}{\linewidth}{llX}
    \toprule
    \textbf{Library identifier} & \textbf{Display name} & \textbf{Description} \\
    \midrule
    H5P.MultiChoice 1.16         & Multiple Choice    & Single- or multi-answer question with up to four options. \\
    H5P.TrueFalse 1.6            & True/False         & Binary true-or-false assertion with optional feedback. \\
    H5P.Blanks 1.14              & Fill in the Blanks & Cloze test with asterisk-delimited correct answers. \\
    H5P.ImageHotspotQuestion 1.8 & Image Hotspot      & Click the correct region on a provided image. \\
    H5P.DragQuestion 1.14        & Drag and Drop      & Drag labelled items to their correct target zones. \\
    H5P.MarkTheWords 1.9         & Mark the Words     & Highlight key terms in a provided text passage. \\
    \bottomrule
  \end{tabularx}
\end{table}

\subsection{Server-Side H5P Libraries}

Traditional H5P deployments embed content within a CMS (WordPress or Moodle)
that handles library management, content storage, and rendering. Lumi
Education GmbH publishes a set of official Node.js libraries---
\texttt{@lumieducation/h5p-server} and \texttt{@lumieducation/h5p-express}
---that expose H5P's core capabilities as a standalone Node.js service,
decoupled from any CMS~\cite{lumieducation2023}. This library provides
H5P library management, content CRUD operations, package generation, and a
rendering API that produces HTML suitable for embedding. The present project
is built on this library, and it is the technical enabler for operating
H5P without WordPress or Moodle.

\subsection{Limitations of WordPress-Based H5P Deployments}

The dominant H5P deployment in Vietnamese universities uses the WordPress
H5P plugin. While functionally complete, this approach imposes several
structural disadvantages at institutional scale:

\begin{itemize}
  \item \textbf{Multi-system complexity.} Instructors must manage a
  WordPress installation (user accounts, posts, media library) in addition
  to H5P content. The cognitive overhead of understanding two distinct
  systems creates unnecessary extraneous load~\cite{sweller1988}.

  \item \textbf{Performance overhead.} WordPress serves every page through a
  full PHP execution stack with database queries, regardless of content type.
  Measured response times on institutional WordPress deployments range from
  2--5 seconds per page, compared with sub-second REST API responses in
  purpose-built single-page applications.

  \item \textbf{Version management friction.} H5P library versions and
  WordPress versions must be kept in sync manually; mismatches produce
  silent failures that are difficult to diagnose without technical expertise.

  \item \textbf{No AI capabilities.} The WordPress H5P plugin has no
  mechanism for analysing video content and suggesting interactive elements.
  Every question must be authored manually from scratch.
\end{itemize}

% ─────────────────────────────────────────────────────────────────────────────
\section{Interactive Video: Research Evidence and Design Principles}

\subsection{Learning Outcome Evidence}

The educational literature provides robust evidence that interactive video
outperforms passive video for learning retention and transfer. Zhang et al.\
conducted one of the earliest controlled experiments comparing interactive
and non-interactive video in an e-learning context and found that the
interactive group scored 20--25\% higher on post-tests and reported
significantly higher satisfaction~\cite{zhang2006}. The mechanism proposed
was that interactive elements force learners to actively process content at
regular intervals, preventing the superficial passive exposure associated
with lecture-video formats.

Merkt et al.\ extended this finding by examining the role of control
features in video-based learning~\cite{merkt2011}. Videos with embedded
questions produced better knowledge retention than equivalent print materials
when learners actively engaged with the interactive features. Critically,
this advantage was most pronounced for university-level content, suggesting
that H5P interactive video is particularly well-matched to the VNU-HCM
context.

Freeman et al.\ produced a meta-analysis across 225 published studies
demonstrating that active learning interventions---of which interactive video
is one class---increased examination scores by an average of 6\% and halved
the failure rate compared to passive lecture formats~\cite{freeman2014}.

\subsection{Embedded Question Design}

Vural identified design principles that distinguish effective from
ineffective interactive video~\cite{vural2013}. Questions should be
positioned \emph{after} a concept has been fully explained, not
mid-sentence. Questions should be spaced at least 30~seconds apart to
avoid fragmenting the learning experience. Each question should test the
single most important idea in the preceding segment. These principles are
directly encoded in the system prompt used by the AI pipeline of this
project (see Chapter~4, Section~4.5.2).

Mayer's Cognitive Theory of Multimedia Learning provides the theoretical
basis for these recommendations~\cite{mayer2009}. Effective multimedia
instruction must manage three cognitive channels---visual, auditory, and
their integration in working memory---and well-designed interactive
elements reduce extraneous load while increasing germane load. A question
that tests exactly what has just been explained supports the integration
channel without taxing working memory with irrelevant processing.

% ─────────────────────────────────────────────────────────────────────────────
\section{User Experience Design and Usability}

\subsection{User-Centred Design}

User-centred design (UCD) is a design philosophy and methodology that
prioritises the needs, goals, and constraints of end users at every stage of
the design process. The ISO~9241-11 standard defines usability as the
extent to which a product can be used by specified users to achieve
specified goals with effectiveness, efficiency, and satisfaction in a
specified context of use~\cite{iso9241}. In educational technology, this
definition must be contextualised to the specific goals and context of
instructors: producing high-quality learning content quickly and without
frustration, within the institutional constraints of time-pressured
academic schedules.

Garrett's layered model of user experience decomposes a product into five
planes: Strategy (objectives), Scope (features), Structure (interaction
design), Skeleton (interface design), and Surface (visual
design)~\cite{garrett2010}. Norman's affordance theory provides a
complementary framework: good interface design makes the correct action
obvious and makes errors difficult~\cite{norman1988}. Both frameworks were
applied during the interface design process described in Chapter~3.

\subsection{Usability Heuristics}

Nielsen's ten usability heuristics provide a practical checklist for
evaluating interface designs against established
principles~\cite{nielsen1994}. The most relevant heuristics for an
educational authoring tool are: (1)~\emph{visibility of system status}---
users must know whether their video is processing, whether AI analysis is
running, and whether a question has been saved; (2)~\emph{match between
system and the real world}---all terminology must correspond to pedagogical
concepts, not technical ones; and (3)~\emph{error prevention}---the system
should make it impossible to submit a malformed H5P question. Molich and
Nielsen demonstrated that applying heuristic evaluation to interface designs
before user testing produces a measurable improvement in
usability~\cite{molich1990}.

\subsection{Measuring Usability: The System Usability Scale}

The System Usability Scale (SUS) is a validated, ten-item questionnaire
that produces a single 0--100 usability score~\cite{brooke1996}. It is
widely used in both academic research and industry practice because of its
simplicity and comparability across studies. A SUS score above 68 is
considered above average; scores above 80 are considered excellent. The
SUS was used as the primary standardised usability instrument in the
evaluation described in Chapter~5.

% ─────────────────────────────────────────────────────────────────────────────
\section{Cognitive Load Theory}

Sweller's cognitive load theory (CLT) identifies three categories of
cognitive load that an interface imposes on a user during a
task~\cite{sweller1988}:

\begin{itemize}
  \item \textbf{Intrinsic load} arises from the inherent complexity of the
  task (e.g., formulating a pedagogically sound multiple-choice question).
  This load is irreducible; it is the productive work the instructor came
  to do.

  \item \textbf{Extraneous load} arises from the design of the interface
  rather than the task itself (e.g., navigating through six screens to reach
  the timestamp entry field). This load is wasteful and should be minimised
  through good design~\cite{chandler1991}.

  \item \textbf{Germane load} arises from the cognitive effort of building
  new knowledge schemas---in this context, the instructor's understanding of
  how their video content maps to interactive learning moments.
\end{itemize}

Chandler and Sweller applied CLT specifically to the format of instructional
materials and found that splitting related information across multiple
screens (the split-attention effect) significantly increased extraneous load
and reduced task accuracy~\cite{chandler1991}. The design implication is
direct: all information needed to author a single H5P interaction---the
video player, the timestamp indicator, the question type selector, and the
question editor---must be co-located on a single screen. Paas, Renkl, and
Sweller confirmed that this integrated format principle consistently reduces
task completion time by 30--50\% compared to split-screen
presentations~\cite{paas2003}.

% ─────────────────────────────────────────────────────────────────────────────
\section{Artificial Intelligence for Educational Content Generation}

\subsection{Large Language Models and Few-Shot Prompting}

The publication of GPT-3 by Brown et al.\ demonstrated that sufficiently
large language models acquire extensive world knowledge during pre-training
and can perform a wide range of tasks from a small number of in-context
examples (few-shot prompting)~\cite{brown2020}. For educational question
generation, this means that a well-constructed prompt can instruct an LLM
to analyse a transcript segment, identify key concepts, and generate
assessment questions that conform to a specified schema---without any
task-specific fine-tuning. Anthropic's Claude family of models builds on
this capability with enhanced instruction-following, constitutional AI
safety training, and more reliable structured JSON output, making it
particularly suitable for tasks requiring schema-constrained
generation~\cite{anthropic2024}.

\subsection{Prompt Engineering for Structured Output}

Wei et al.\ introduced chain-of-thought prompting, showing that instructing
a model to reason step-by-step before producing an answer substantially
improves accuracy on complex tasks~\cite{wei2022}. For question generation,
this approach is realised as a prompting strategy that asks the model to
first identify the key teaching moment in a transcript segment, then assign
the appropriate Bloom's taxonomy level, and finally generate the question
and answer options---rather than generating the final output directly.

White et al.\ catalogued a set of prompt patterns that improve the
reliability of structured LLM output~\cite{white2023}. The most relevant
to this project is the \emph{output automater} pattern, which instructs the
model to produce its response as a JSON array conforming to a precisely
specified schema. This pattern, combined with runtime validation via the
Zod TypeScript schema library, forms the core of the AI pipeline's output
guarantees (detailed in Chapter~4, Section~4.5).

\subsection{Bloom's Revised Taxonomy for Question Classification}

Bloom's taxonomy provides a hierarchical classification of cognitive
objectives, ranging from basic recall (\emph{remember}) to analysis and
synthesis (\emph{create})~\cite{bloom1956}. Anderson and Krathwohl's
revised taxonomy updated the original framework with active verb labels and
clarified the distinction between lower-order (remember, understand, apply)
and higher-order (analyse, evaluate, create) cognitive
skills~\cite{anderson2001}. The AI pipeline assigns a Bloom's level to
each generated suggestion, enabling instructors to ensure that their
interactive video includes a deliberate spread of cognitive difficulty
without requiring them to possess specialist instructional design knowledge.

\subsection{Automatic Speech Recognition}

Radford et al.\ introduced OpenAI Whisper, a large-vocabulary ASR model
trained on 680,000 hours of multilingual audio using weak
supervision~\cite{radford2023}. Whisper achieves near-human word error
rates on English audio and competitive performance on Vietnamese. The
\texttt{faster-whisper} implementation is used in this project for on-device
audio transcription of uploaded video files where no caption file is
available. The quality of transcription directly affects the quality of
AI-generated suggestions, motivating the choice of Whisper over simpler
ASR approaches.

% ─────────────────────────────────────────────────────────────────────────────
\section{Cloud-Native Web Architecture}

\subsection{Three-Tier Architecture}

The three-tier architecture---comprising a client (presentation) tier, an
application (logic) tier, and a data tier---is the standard pattern for
web-based applications requiring clear separation of concerns and
independent scalability of components~\cite{bass2012}. Armbrust et al.\
demonstrated that this architecture achieves a 40--60\% reduction in total
cost of ownership compared to monolithic server deployments while improving
availability through redundancy~\cite{armbrust2010}. For this project, the
model maps to: the React~18 SPA (client), the Express.js REST API server
(application), and the SQLite database (data).

\subsection{Single-Page Applications with React}

A Single-Page Application (SPA) delivers a complete JavaScript application
on the initial page load and subsequently manages all rendering and
navigation without full page reloads~\cite{react2023}. SPAs are particularly
well suited to interactive authoring tools because they eliminate the latency
and state-loss associated with traditional multi-page web applications.
React~18's concurrent rendering features allow the browser to remain
responsive during long-running operations such as AI analysis, which can
take 5--15 seconds, while streaming token updates are progressively displayed.

TypeScript adds static type checking to JavaScript, catching type mismatches
at compile time rather than at runtime~\cite{typescript2023}. For a project
whose AI pipeline must map between several different JSON schemas (transcript
segments, AI suggestions, H5P content parameters), TypeScript's type system
provides a critical layer of correctness assurance.

\subsection{Real-Time Communication with Server-Sent Events}

Server-Sent Events (SSE) is a W3C standard for unidirectional,
server-to-client streaming over a persistent HTTP connection. Unlike
WebSockets, SSE requires no connection upgrade and works through standard
HTTP proxies. The AI analysis streaming endpoint in this project uses SSE
to push token-level progress updates from the Claude or Ollama inference
session to the browser, giving instructors immediate feedback that analysis
is progressing rather than waiting without feedback for up to 15 seconds.

% ─────────────────────────────────────────────────────────────────────────────
\section{LMS Integration and Interoperability Standards}

\subsection{Learning Tools Interoperability}

The Learning Tools Interoperability (LTI) specification, maintained by IMS
Global Learning Consortium, defines a standard protocol for integrating
external tools into Learning Management Systems~\cite{lti2019}. LTI~1.3
uses OAuth~2.0 for authentication and provides mechanisms for tool
launches, grade passback, and deep linking to specific content items.
Severance et al.\ found that LTI-compliant tools achieved higher adoption
rates than custom-integrated tools because they could be configured once
and used across all LTI-compatible LMSs without modification~\cite{severance2010}.
The H5P packages exported by this platform are directly importable into any
LTI-compatible LMS; the backend additionally exposes LTI launch endpoints
for direct integration.

\subsection{Learning Analytics}

Siemens and Long argued that the shift from passive to active digital
learning creates a data exhaust that, properly analysed, can improve
instruction quality and learner outcomes~\cite{siemens2011}. Interactive
video creates particularly rich interaction data: every question attempt,
timestamp seek, replay event, and export action is a potential analytics
data point. The current platform captures basic analytics (video visit
counts, export events) for administrative review. Comprehensive learning
analytics integration with the VNU-HCM MOOC platform is identified as a
priority direction for future work.

% ─────────────────────────────────────────────────────────────────────────────
\section{Related Work}

\subsection{WordPress H5P Plugin}

The official WordPress H5P plugin provides full-featured H5P authoring
within a WordPress CMS~\cite{h5p2023}. As discussed in Section~2.2.3, this
approach imposes significant usability overheads. It has no AI capabilities
and provides no real-time preview within the authoring interface. It remains
the dominant H5P authoring environment in Vietnamese universities.

\subsection{Lumi H5P Desktop Editor}

Lumi, developed by Lumi Education GmbH, is a cross-platform desktop
application providing a standalone H5P editor without requiring a
CMS~\cite{lumieducation2023}. It eliminates the WordPress dependency and
offers a cleaner authoring experience for individual instructors. However,
it is a desktop application (not a web service), does not support
collaborative or centralised authoring, and has no AI integration.

\subsection{Moodle H5P Integration}

Moodle's built-in H5P activity module provides another authoring path for
H5P content~\cite{moodle2023}. While tightly integrated with the LMS
gradebook and user management, the Moodle H5P editor shares many of the
usability limitations of the WordPress plugin and provides no AI assistance.

\subsection{Commercial Authoring Platforms}

Articulate Storyline and Adobe Captivate offer sophisticated interactive
video authoring but carry licensing costs (typically \$1{,}000--\$2{,}000
USD per seat per year) that are prohibitive for Vietnamese public
universities. Neither produces native H5P output, limiting LMS
compatibility. Neither has integrated AI question generation.

% ─────────────────────────────────────────────────────────────────────────────
\section{Research Gap}

The review above establishes a clear and multi-dimensional gap in the
existing literature and available tools. No current system simultaneously:

\begin{enumerate}
  \item provides a \emph{web-based}, CMS-free H5P interactive video
  authoring environment suitable for institutional deployment;
  \item integrates a \emph{transcript-to-H5P AI pipeline} using a
  production-grade LLM with structured output validation;
  \item applies \emph{Bloom's taxonomy classification} to AI-generated
  suggestions to guide pedagogical design;
  \item has been \emph{empirically evaluated} against a real-world authoring
  baseline and reported with quantitative usability metrics.
\end{enumerate}

This project addresses all four gaps simultaneously, making a contribution
that is both technically novel and practically significant for the VNU-HCM
educational ecosystem and for the broader field of AI-assisted educational
content authoring.
